\chapter{Quantum corrections to Arnold relations}
\label{ch:quantum-corrections}

The universal modular Maurer--Cartan element
$\Theta_\cA \in \MC(\mathfrak g^{\mathrm{mod}}_\cA)$ satisfies
\[
D_\cA\Theta_\cA
  + \tfrac{1}{2}[\Theta_\cA,\Theta_\cA]=0 .
\]
The factor $\tfrac{1}{2}$ is part of the convention: it is the same
factor as in the BV quantum master equation and in the operadic
Maurer--Cartan equation of Appendix~\ref{app:signs}.  The
genus-$0$ collision component is the Arnold flatness identity for the
logarithmic forms $\eta_{ij}=d\log(z_i-z_j)$.  At positive genus the
total modular equation remains a flat MC equation; the new term is the
fiberwise Hodge curvature.  Its scalar projection is
$\kappa(\cA)\lambda_g$ on the uniform-weight scalar lane, while
interacting multi-weight algebras carry additional cross-channel
corrections invisible to that scalar projection
(Chapter~\ref{ch:genus-expansions}).

\begin{definition}[Quantum corrections]\label{def:quantum-corrections}
\index{quantum corrections|textbf}
The \emph{quantum corrections} of~$\cA$ are the positive-genus
components in the completed stable-graph expansion
\[
 D = D^{(0)}+\sum_{g\geq 1}\hbar^g D^{(g)} .
\]
The genus-$0$ part $D^{(0)}$ is the chiral bar differential on
configuration spaces of~$X$, governed by the Arnold relation for
$\eta_{ij}=d\log(z_i-z_j)$ on $\FM_n(X)$.  The terms $D^{(g)}$ for
$g\geq 1$ are the modular corrections: they record non-separating
clutching, period dependence, and the curvature component of the
completed convolution algebra.  The formal variable~$\hbar$ has
degree~$0$ and records only the genus filtration.
\end{definition}

\begin{remark}[Genus-filtration parameter]
\label{rem:qc-hbar-genus-filtration}
The symbol $\hbar$ in Definition~\ref{def:quantum-corrections} is the
Rees parameter for the genus filtration of the modular convolution
algebra.  A stable graph of genus
$g=h^1(\Gamma)+\sum_{v}g_v$ lies in the coefficient of~$\hbar^g$;
one application of the non-separating clutching operator~$\Delta$
raises this exponent by one.  The physical partition-function
normalization uses $\hbar_{\mathrm{phys}}^{2g-2}$ and the scalar
free-energy convention uses the even variable
$\hbar_{\mathrm{phys}}^2$.  The imported K3 boundary parameters in
\S\ref{sec:qc-supplement} are Vol.~III quantum-group deformation
parameters.  These three parameters live in separate filtrations.
\end{remark}

The genus-$0$ identity
$\eta_{12} \wedge \eta_{23} + \eta_{23} \wedge \eta_{31}
+ \eta_{31} \wedge \eta_{12}=0$ is the genus-$0$ projection of the
chain-level equation in $\Convinf(\barBch(\cA),\cA)$.  In genus
$g\geq1$ the completed stable-graph sum splits after choosing the
strong-generator channel set $\mathrm{Ch}(\cA)$:
\[
F_g(\cA)
=
\sum_{\substack{\Gamma,\ \sigma:E(\Gamma)\to\mathrm{Ch}(\cA)\\
                \sigma\ \mathrm{constant}}}
\frac{W_\Gamma(\sigma)}{|\operatorname{Aut}(\Gamma)|}
\;+\;
\sum_{\substack{\Gamma,\ \sigma:E(\Gamma)\to\mathrm{Ch}(\cA)\\
                \sigma\ \mathrm{not\ constant}}}
\frac{W_\Gamma(\sigma)}{|\operatorname{Aut}(\Gamma)|}
=F_g^{\mathrm{diag}}(\cA)+\delta F_g^{\mathrm{cross}}(\cA).
\]
The Chern--Weil map from the diagonal summand to the Hodge lane sends
$F_g^{\mathrm{diag}}(\cA)$ to
$\kappa(\cA)\lambda_g^{\mathrm{FP}}$; the obstruction-class version is
$\operatorname{obs}_g^{\mathrm{scal}}(\cA)=\kappa(\cA)\lambda_g$.
Heisenberg has only this scalar curvature channel.  Affine
Kac--Moody algebras add the cubic Lie channel.  The $\beta\gamma$
system has the free-field mixed-channel cancellation recorded in
Chapter~\ref{ch:genus-expansions}.  Principal $\mathcal W_N$ algebras
for $N\geq 3$ are the standard interacting multi-weight witnesses:
their diagonal scalar term is $\kappa(\cA)\lambda_g$, and their full
free energy contains $\delta F_g^{\mathrm{cross}}$ for $g\geq 2$.

\section{Braid groups and Arnold relations}

\subsection{Arnold relations and the braid group}

Arnold's relation is the codimension-two collision identity behind
flatness of the logarithmic configuration-space connection.  For
three ordered points on the affine line, write
$z_{ij}=z_i-z_j$ and
\[
\eta_{ij}=d\log z_{ij}=\frac{dz_i-dz_j}{z_i-z_j}.
\]
The forms $\eta_{ij}$ are the de~Rham representatives of the
Orlik--Solomon generators for the braid arrangement; the braid-group
monodromy is obtained by integrating them around the collision
divisors.

\subsubsection{The braid derivation of Arnold relations}

The local winding normalization is
\[
\frac{1}{2\pi i}\oint_{|z_{ij}|=\epsilon}\eta_{ij}=1,
\]
while the quadratic relation follows from the partial-fraction
identity for the three collision hyperplanes.  Explicitly,
\begin{align}
\eta_{12} &= d\log(z_1 - z_2) = \frac{dz_1 - dz_2}{z_1 - z_2} \\
\eta_{23} &= d\log(z_2 - z_3) = \frac{dz_2 - dz_3}{z_2 - z_3} \\
\eta_{31} &= d\log(z_3 - z_1) = \frac{dz_3 - dz_1}{z_3 - z_1}
\end{align}
and clearing denominators in
\[
\frac{1}{z_{12}z_{23}}
+\frac{1}{z_{23}z_{31}}
+\frac{1}{z_{31}z_{12}}=0
\]
gives $z_{31}+z_{12}+z_{23}=0$.  Multiplying by the common
two-form numerator gives Arnold's relation
\[
\eta_{12}\wedge\eta_{23}
+\eta_{23}\wedge\eta_{31}
+\eta_{31}\wedge\eta_{12}=0.
\]
This is the form-level relation used in the bar differential; the
Orlik--Solomon sign convention for generators is compared with the
differential-form convention in Appendix~\ref{app:signs}.

\subsection{The meaning of integrability}

These relations are the integrability conditions for the genus-$0$
geometric bar complex.  On the affine Kac--Moody comparison surface
they are the flatness equations for the Knizhnik--Zamolodchikov
connection on conformal blocks \cite[Section~3.4]{BD04}; for a
general chiral algebra they are the codimension-two cancellation in
the residue differential.

Consider the bar complex:
\[\bar{B}^{\text{geom}}(\mathcal{A}) = \bigoplus_n \Gamma(\overline{C}_n(X), \mathcal{A}^{\boxtimes n} \otimes \Omega^*_{\log})\]

with differential $d = d_{\text{internal}} + d_{\text{residue}} + d_{\text{deRham}}$. The condition $d^2 = 0$ says that the differential defines a flat connection on an infinite-dimensional bundle.

\subsubsection{The Maurer--Cartan perspective}

More precisely, we can view $d$ as a connection on the graded vector bundle:
\[\mathcal{E} = \bigoplus_{n,k} \mathcal{A}^{\boxtimes n} \otimes \Omega^k_{\log}\]

The flatness condition is the Maurer--Cartan equation
\[d\omega + \frac{1}{2}[\omega, \omega] = 0\]

where $\omega = \sum_{i<j} t_{ij}\, \eta_{ij}$ is the residue-valued
logarithmic coefficient in the geometric bar differential, with
$t_{ij} \in \mathrm{End}(\mathcal{A}^{\otimes n})$ acting on the $i$-th
and $j$-th tensor factors via the chiral bracket and
$\eta_{ij} = d\log(z_i - z_j)$ the local collision form. On an affine
genus-$0$ screen the Arnold relations are precisely the conditions
ensuring this equation holds; on $\mathbb{P}^1$ and on higher-genus
curves one must add the projective or mixed period relations of
Proposition~\ref{prop:arnold-higher-genus}.

\subsubsection{Concrete computation}

For $n=3$, the de~Rham squares $d\eta_{ij}=0$ do not carry the
constraint.  The non-trivial term in $d^2$ is the sum of iterated
residues at the triple collision divisor.  Its form factor is exactly
\[
\eta_{12}\wedge\eta_{23}
+\eta_{23}\wedge\eta_{31}
+\eta_{31}\wedge\eta_{12},
\]
so the Arnold relation is the local cancellation that removes the
codimension-two residue.

\section{Genus one: theta functions and central extensions}

\subsection{The genus one quantum correction}

On the torus $E_\tau=\mathbb C/(\mathbb Z+\tau\mathbb Z)$ there is no
global holomorphic function with one simple zero and strict double
periodicity.  The genus-$1$ propagator is therefore computed in a
holomorphic trivialization on the universal cover and then projected to
the invariant scalar curvature class.

\subsubsection{The prime-form construction}

The genus-$1$ prime form in the standard holomorphic
trivialization is
\[
E(z|\tau)=\frac{\theta_1(z|\tau)}{\theta_1'(0|\tau)}.
\]
Its logarithmic derivative
\[
K_\tau(z)=\partial_z\log E(z|\tau)
\]
is the propagator used in the bar differential.  It is related to the
Weierstrass sigma function by
\[
K_\tau(z)=\zeta_\sigma(z|\tau)-2\eta_1 z,
\qquad
\zeta_\sigma(z|\tau)=\partial_z\log\sigma(z|\tau).
\]
This distinction is essential: the plain sigma zeta function has no
$G_2z$ term in its Laurent expansion, while the prime-form derivative
does.  The genus-$1$ scalar curvature below is computed from
$K_\tau$, not from $\zeta_\sigma$ alone.

\subsubsection{The quasi-periodicity and its consequences}

The sigma zeta function is quasi-periodic:
\begin{align}
\zeta_\sigma(z + 1|\tau) &= \zeta_\sigma(z|\tau) + 2\eta_1 \\
\zeta_\sigma(z + \tau|\tau) &= \zeta_\sigma(z|\tau) + 2\eta_\tau
\end{align}

where the quasi-periods satisfy the Legendre relation
\[\tau\eta_1 - \eta_\tau = \pi i\]
from the residue normalization
$\oint_{\partial\mathcal F}\zeta_\sigma(z)\,dz=2\pi i$.  Hence
$K_\tau(z+1)=K_\tau(z)$ and
$K_\tau(z+\tau)=K_\tau(z)-2\pi i$.  The $B$-cycle shift is the
genus-$1$ source of the Hodge curvature.

\subsubsection{Computing the quantum correction}

In the prime-form trivialization set
\[
\eta_{ij}^{(1)}
 = d\log E(z_i-z_j|\tau)
 = K_\tau(z_i-z_j)(dz_i-dz_j).
\]
The Arnold combination is
\[\mathcal{A}_3^{(1)} = \eta_{12}^{(1)} \wedge \eta_{23}^{(1)} + \eta_{23}^{(1)} \wedge \eta_{31}^{(1)} + \eta_{31}^{(1)} \wedge \eta_{12}^{(1)}\]

We derive the scalar defect from the Eisenstein-regularised Laurent
expansion of~$K_\tau$.

\emph{Step~1: Form-factor reduction.}
Write $u = z_{12}$, $v = z_{23}$, $w = z_{31} = -(u+v)$
and $\alpha = dz_1 - dz_2$, $\beta = dz_2 - dz_3$.
Note that $dz_3 - dz_1 = -(\alpha + \beta)$, so
\[
 \beta \wedge (-\alpha - \beta) = -\beta \wedge \alpha = \alpha \wedge \beta,
 \qquad
 (-\alpha - \beta) \wedge \alpha = -\beta \wedge \alpha = \alpha \wedge \beta.
\]
Every cyclic term carries the \emph{same} form factor, giving
\begin{equation}\label{eq:arnold-genus1-factored}
\mathcal{A}_3^{(1)}
= \bigl[K_\tau(u)K_\tau(v)
 + K_\tau(v)K_\tau(w) + K_\tau(w)K_\tau(u)\bigr]
 \cdot \alpha \wedge \beta.
\end{equation}

\emph{Step~2: Pole cancellation.}
Denote the coefficient function
$S(u,v)=\sum_{\mathrm{cyc}}K_\tau(u)K_\tau(v)$.
At $z_1 = z_2$ (i.e.\ $u = 0$), the first and third terms
each have a simple pole in~$u$, with residues $K_\tau(v)$ and
$K_\tau(w) = K_\tau(-v) = -K_\tau(v)$ respectively; the second
term is regular. So the residue of~$S$ at $u = 0$ vanishes.
By symmetry the same holds at $v = 0$ and $w = 0$.
Thus the collision-pole part vanishes; only the regular period
polynomial of the chosen trivialization can contribute to the scalar
projection.

\emph{Step~3: Laurent expansion via Eisenstein summation.}
The prime-form derivative admits the Eisenstein-regularised expansion
(cf.~\S\ref{sec:km-bar-complex}):
\begin{equation}\label{eq:primeform-eisenstein-expansion}
K_\tau(z)
= \frac{1}{z} - G_2(\tau)\,z - G_4(\tau)\,z^3 - G_6(\tau)\,z^5 - \cdots
\end{equation}
where $G_{2k}(\tau) = \sum'_{(m,n)} (m + n\tau)^{-2k}$
and $G_2(\tau) = \tfrac{\pi^2}{3}\,E_2(\tau)$ requires
Eisenstein summation (inner sum over~$m$ first).  Equivalently,
$2\eta_1=G_2(\tau)$ in the sigma normalization and
$K_\tau=\zeta_\sigma-G_2z$.

\emph{Step~4: Constant term of $S$.}
Substituting~\eqref{eq:primeform-eisenstein-expansion} into $S$
and collecting powers of $(u,v)$, the singular terms cancel
by Step~2. The leading contributions are:
\begin{align}
\text{from }1/(uv):\quad
 &\textstyle\sum_{\mathrm{cyc}} \frac{1}{uv}
 = \frac{u + v + w}{uvw} = 0,
 \label{eq:arnold-leading}\\[4pt]
\text{from }G_2:\quad
 &\textstyle -G_2 \sum_{\mathrm{cyc}}\bigl(\tfrac{u}{v}
 + \tfrac{v}{u}\bigr)
 = -G_2 \cdot \frac{-(u^3 + v^3 + w^3)}{uvw}
 = -G_2 \cdot \frac{-3uvw}{uvw} = 3G_2,
 \label{eq:arnold-G2}\\[4pt]
\text{from }G_4, G_2^2:\quad
 &(5G_4 - G_2^2)(u^2 + uv + v^2) + \cdots
 \label{eq:arnold-higher}
\end{align}
In~\eqref{eq:arnold-G2} we used Newton's identity
$u^3 + v^3 + w^3 = 3uvw$ (valid when $u + v + w = 0$), and
in~\eqref{eq:arnold-leading} the genus-$0$ Arnold relation.

\emph{Step~5: Scalar projection.}
The position-dependent terms in~\eqref{eq:arnold-higher} are not
identified here with zero in $H^{2,0}(C_3(E_\tau))$: exactness on the
universal cover does not by itself descend to the torus configuration
space, and the prime-form derivative carries $B$-cycle
quasi-periodicity.  The scalar curvature lane keeps only the
translation-invariant constant term of the Eisenstein expansion.
Denote this projection by $\pi_{\mathrm{scal}}$.

\emph{Step~6: The scalar Arnold defect.}
Combining Steps~4 and~5 gives
\begin{equation}\label{eq:arnold-defect-genus1}
\boxed{
\pi_{\mathrm{scal}}\mathcal{A}_3^{(1)}
= 3\,G_2(\tau) \cdot \alpha \wedge \beta
= \pi^2 E_2(\tau) \cdot \alpha \wedge \beta
}
\end{equation}
where $E_2(\tau) = 1 - 24\sum_{n=1}^\infty \sigma_1(n)\,q^n$
is the weight-$2$ quasi-modular Eisenstein series
($q = e^{2\pi i \tau}$).
The factor $\pi^2$ belongs to the analytic generator
$G_2=(\pi^2/3)E_2$ and is removed by the scalar-lane normalization;
the invariant $\kappa(\cA)$ is fixed by the degree-two OPE
coefficient.  Put
\[
e_2^{\mathrm{an}}
:=\pi^{-2}\pi_{\mathrm{scal}}\mathcal{A}_3^{(1)}
=E_2(\tau)\,\alpha\wedge\beta .
\]
The analytic-to-Hodge normalization
\[
\nu_1\colon \mathbb C\cdot e_2^{\mathrm{an}}
\longrightarrow H^2(\overline{\mathcal M}_{1,1};\mathbb C),
\qquad
\nu_1(e_2^{\mathrm{an}})=\lambda_1
\]
is the scalar lane used by Theorem~\ref{thm:genus-universality}.
Contracting the degree-two OPE coefficient with this generator gives
the genus-$1$ scalar curvature in the $E_2$ trivialization,
\[
m_{0,\mathrm{scal}}^{(1)}(\cA)=\kappa(\cA)\,E_2(\tau),
\]
equivalently $F_1(\cA)=\kappa(\cA)/24$ after the Hodge projection and
integration over $\overline{\mathcal M}_{1,1}$.
The non-holomorphic completion $\hat{E}_2(\tau) = E_2(\tau)
- 3/(\pi\,\mathrm{Im}(\tau))$ governs the single-valued
propagator and produces the Arakelov $(1,1)$-form
$2\pi i \cdot \omega_\tau$ of
Theorem~\ref{thm:arnold-higher-genus}(b).

\begin{remark}[Scalar one-loop class and determinant amplitude]
\label{rem:qc-scalar-one-loop-vs-determinant}
The equality $F_1(\cA)=\kappa(\cA)/24$ is the Hodge intersection
number $\kappa(\cA)\int_{\overline{\mathcal M}_{1,1}}\lambda_1$.
The full torus determinant is a separate physical amplitude.  For an
even lattice VOA of rank~$d$, the physical genus-$1$ partition
function is
\[
 Z_1(V_\Lambda,\tau)=\Theta_\Lambda(\tau)\,\eta(\tau)^{-d},
 \qquad
 \partial_\tau\log\eta(\tau)=\frac{\pi i}{12}E_2(\tau),
\]
so the oscillator determinant contains the whole
$q$-series of~$\eta^{-d}$, while the scalar lane retains only the
Hodge coefficient $d/24$.  For the free $\beta\gamma$ system the
determinantal statement is a statement about
$\det R\Gamma(\Sigma_g,K^\lambda)$ and its Quillen refinement; this
chapter uses only the induced Mumford/Hodge scalar unless a separate
determinant-line comparison theorem is invoked.  The scalar
obstruction formula therefore uses $\kappa(\cA)/24$, with
$-\log\eta(\tau)$ and zeta-regularized determinants remaining in the
determinant-line lane.
\end{remark}

\subsection{The central extension at genus one}

\subsubsection{The central extension}

Let $\mathcal H_k$ be the rank-one Heisenberg vertex algebra at level
$k$, generated by modes $a_n$ with
\[
[a_m,a_n]=k\,m\,\delta_{m+n,0}\mathbf 1 .
\]
This central extension is a genus-$0$ OPE datum.  Genus~$1$ leaves
the commutation relations fixed and couples the same scalar level
$k=\kappa(\mathcal H_k)$ to the Hodge curvature represented above by
$E_2(\tau)$.

\subsubsection{The explicit construction of the central element}

For the Heisenberg algebra the scalar genus-$1$ curvature is
\[
m_{0,\mathrm{scal}}^{(1)}(\mathcal H_k)=k\,E_2(\tau),
\qquad
F_1(\mathcal H_k)=k/24 .
\]

The identity
$\partial_\tau\log\eta(\tau)=\frac{\pi i}{12}E_2(\tau)$ explains the
same $E_2$ coefficient in the genus-$1$ character.  The central charge
$c=1$ for a rank-one Heisenberg algebra is a separate genus-$0$ VOA
datum; the modular characteristic used here is the level
$\kappa(\mathcal H_k)=k$.

\subsubsection{The cocycle condition}

The central extension is classified by the standard cocycle on the
abelian loop algebra:
\[
\omega_k(a_m,a_n)=k\,m\delta_{m+n,0}.
\]

Since the underlying loop algebra is abelian, the cocycle condition is
immediate:
\[\omega([a_\ell, a_m], a_n) + \omega([a_m, a_n], a_\ell) + \omega([a_n, a_\ell], a_m) = 0\]

For non-abelian affine algebras the analogous cocycle uses the
trace-form pairing and must be kept separate from the KZ normalization;
see Chapter~\ref{chap:kac-moody} and
\eqref{eq:qc-trace-kz-normalizations} below.

\subsection{Trace-form and KZ normalizations}

For an affine Kac--Moody algebra $V_k(\mathfrak g)$ with normalized
invariant form $\kappa_{\mathfrak g}$,
\[
J^a(z)J^b(w)
\sim
\frac{k\,\kappa_{\mathfrak g}(a,b)}{(z-w)^2}
+\frac{f^{ab}{}_cJ^c(w)}{z-w}.
\]
The double pole is the trace-form curvature input; the KZ connection
uses the Sugawara-shifted inverse form.  The two normalizations are
\begin{equation}\label{eq:qc-trace-kz-normalizations}
r^{\mathrm{coll}}_{\widehat{\mathfrak g}_k}(z)
 = \frac{k\,\Omega_{\mathrm{tr}}}{z},
\qquad
r^{\mathrm{KZ}}_{\widehat{\mathfrak g}_k}(z)
 = \frac{\Omega_{\mathrm{KZ}}}{(k+h^\vee)z}.
\end{equation}
Thus $r^{\mathrm{coll}}$ vanishes at $k=0$, while the KZ coefficient
remains nonzero; at the critical level $k=-h^\vee$ the KZ denominator
degenerates, while the trace-form collision residue has no such
denominator.  The modular characteristic is the separate scalar
\[
\kappa(V_k(\mathfrak g))
 = \frac{\dim(\mathfrak g)(k+h^\vee)}{2h^\vee},
\]
and the genus-curvature formula uses this value separately from the
collision level~$k$ and the KZ coefficient $(k+h^\vee)^{-1}$.

\section{Higher genus}

\subsection{\texorpdfstring{Genus $2$}{Genus 2}}

The moduli space $\mathcal{M}_2$ is three-dimensional and is locally
parametrized by the period matrix
\[\Omega = \begin{pmatrix} \tau_{11} & \tau_{12} \\ \tau_{12} & \tau_{22} \end{pmatrix} \in \mathcal{H}_2\]

living in the Siegel upper half-space, the space of symmetric complex $2 \times 2$ matrices with positive definite imaginary part.

\subsubsection{The theta functions}

There are $16$ theta characteristics in genus~$2$: six odd and ten
even.  The universal prime form is written using an odd
characteristic; fermionic determinants and Pfaffian lines depend on
spin parity, but the scalar bar curvature below is the Hodge-class
projection.  Quantum corrections outside this scalar projection may
retain their characteristic parity.

\subsubsection{Detailed computation of genus 2 corrections}

The prime form may be written as
\[E(z,w) = \frac{\theta[\delta](z-w|\Omega)}{h_\delta(z)^{1/2} h_\delta(w)^{1/2}}\]

where $\delta$ is an odd characteristic and $h_\delta$ is the
associated half-differential.  With the convention of
Appendix~\ref{app:signs}, $E(z,w)$ is a section of
$K^{-1/2}\boxtimes K^{-1/2}$ and is antisymmetric.

The logarithmic forms become:
\[\eta_{ij}^{(2)} = d\log E(z_i, z_j) = \partial_i \log E(z_i, z_j)\, dz_i + \partial_j \log E(z_i, z_j)\, dz_j\]

The Arnold combination on $\Sigma_2$ is
\[\mathcal{A}_3^{(2)} = \sum_{\text{cyclic}} \eta_{ij}^{(2)} \wedge \eta_{jk}^{(2)}.\]
As at genus~$1$, the local pole terms cancel on the scalar diagonal
after projection to the Hodge class.  The universal scalar statement is
the cohomological identity
\[
m_{0,\mathrm{scal}}^{(2)}(\cA)
  = \kappa(\cA)\lambda_2
\]
on the uniform-weight scalar lane (Theorem~\ref{thm:genus-universality}).
For a multi-weight algebra the full genus-$2$ free energy is
\[
F_2(\cA)
=\kappa(\cA)\lambda_2^{\mathrm{FP}}
 +\delta F_2^{\mathrm{cross}}(\cA),
\]
where $\delta F_2^{\mathrm{cross}}$ is the nonconstant channel-coloring
subsum of the genus-$2$ stable graphs
(Theorem~\ref{thm:multi-weight-genus-expansion}).  For the first
interacting witness,
\[
\delta F_2^{\mathrm{cross}}(\cW_3)
=\frac{c+204}{16c}.
\]
The planted-forest contribution lives in the higher-degree shadow
tower and leaves the degree-$2$ scalar projection unchanged
(Remark~\ref{rem:pf-does-not-affect-scalar-lane}).

\section{\texorpdfstring{The $A_\infty$ structure}{The A-infinity structure}}

The $A_\infty$ structure (Stasheff~\cite{Sta63}) encodes the
tree-level homotopies.  Positive genus adds curvature and stable-graph
operations in the completed modular convolution algebra; the
genus-$0$ bar differential remains nilpotent.

\subsection{\texorpdfstring{The complete $A_\infty$ structure}{The complete A-infinity structure}}

An $A_\infty$ algebra consists of operations $m_n: A^{\otimes n} \to A[2-n]$ for $n \geq 1$, satisfying:
\[\sum_{\substack{r+s+t=n \\ r,t \geq 0,\; s \geq 1}} (-1)^{rs+t}\, m_{r+1+t}(\mathrm{id}^{\otimes r} \otimes m_s \otimes \mathrm{id}^{\otimes t}) = 0\]

\subsubsection{For the bar complex}

The transferred bar model carries operations

\[m_1 = d_{\text{bar}}\]

\[m_2(a \otimes b) = \text{Res}_{z_1=z_2}\left[\frac{a(z_1)b(z_2)}{z_1-z_2}\right]\]

\[m_3(a \otimes b \otimes c) = \int_{\overline{C}_3(X)} a(z_1)\, b(z_2)\, c(z_3)\, \eta_{123}^{\text{FM}}\]
where $\eta_{123}^{\text{FM}}$ is the propagator form on the Fulton--MacPherson compactification $\overline{C}_3(X)$, defined via homotopy transfer (cf.\ Appendix~\ref{app:homotopy-transfer}). Concretely, $m_3$ sums over binary trees with 3 leaves, each tree contributing an iterated composition of $m_2$ with the contracting homotopy $h$:
\[m_3(a, b, c) = m_2(m_2(a, b) \cdot h,\, c) + m_2(a,\, h \cdot m_2(b, c))\]

\subsection{Explicit computations for specific algebras}

\subsubsection{For the Heisenberg algebra}

The genus-$0$ Heisenberg transfer is strict: $m_0=0$,
$m_1^2=0$, $m_2$ is the standard product, and $m_n=0$ for
$n\geq 3$.  Positive-genus scalar corrections enter through the
separate modular curvature
$m_{0,\mathrm{scal}}^{(g)}=k\,\lambda_g$ for
$\mathcal H_k$; the Heisenberg tree operations remain strict.

\subsubsection{\texorpdfstring{For the $\beta\gamma$ system}{For the beta-gamma system}}

For a $\beta\gamma$ system with a chosen background charge, the curved
BRST presentation has $m_1^2(a)=[m_0,a]$.  The coefficient of $m_0$
is fixed by the chosen OPE and background-charge convention.  The
ternary operation $m_3$ below is only the genus-$0$ homotopy-transfer
operation; the free-field exactness statement in
Chapter~\ref{ch:genus-expansions} says that the mixed positive-genus
cross-channel graph sum vanishes for this system.

\subsubsection{\texorpdfstring{Explicit computation of $m_3$ for $\beta\gamma$}{Explicit computation of m 3 for beta-gamma}}

The transferred operation $m_3$ is computed by homotopy transfer from
the Fulton--MacPherson compactification $\overline{C}_3(\mathbb{C})$:
\begin{equation}
m_3(a, b, c)
= \sum_{T \in \mathrm{PBT}_3}
 (\pm)\, p \circ \mu_T \circ i^{\otimes 3}(a, b, c),
\end{equation}
where $\mathrm{PBT}_3$ is the set of planar binary trees with $3$
leaves (two trees), each internal edge carrying the contracting
homotopy~$h$, and $p, i$ are the projection and inclusion of the
homotopy retract. This is a genus-$0$ (tree-level) operation;
genus $g \geq 1$ corrections enter through the propagator and the
cohomology of $\overline{\mathcal{M}}_g$; the degree of the
$A_\infty$ operation remains tree-level data.

\subsubsection{For $\mathcal{W}$-algebras}

For principal $\mathcal{W}$-algebras obtained by Drinfeld--Sokolov
reduction, the transferred $A_\infty$ operations $m_n$ are
determined by the BRST spectral sequence: the ghost-number
filtration produces the $\mathcal{W}$-algebra generators on the
$E_2$ page, and the higher $m_n$ encode the
transferred homotopy from the affine Kac--Moody input
(Chapter~\ref{chap:w-algebra-koszul}). The operations $m_n$ are
genus-$0$ (tree-level); their interaction with higher-genus
corrections is governed by the curved $A_\infty$ structure
described in Part~\ref{part:bar-complex}.

\section{Koszul duality and complementary deformations}

\subsection{The fundamental theorem}

\begin{remark}[Koszul complementarity at higher genus]
\label{rem:koszul-complementarity-higher-genus}
\index{quantum corrections!complementarity formula}
The proved complementarity statement used by this chapter is scalar:
on the scalar lane,
\[
F_g^{\mathrm{diag}}(\cA)+F_g^{\mathrm{diag}}(\cA^!)
  =\bigl(\kappa(\cA)+\kappa(\cA^!)\bigr)\lambda_g^{\mathrm{FP}}.
\]
A direct-sum formula for full quantum-correction spaces requires the
extra hypotheses in Theorem~\ref{thm:quantum-complementarity-main};
the scalar curvature identity alone gives only the displayed diagonal
statement.
\end{remark}

\subsection{Examples of Koszul complementarity}

\subsubsection{Example 1: free fermions and free bosons}

The free fermion system $\mathcal{F}$ with OPE:
\[\psi(z)\psi(w) \sim \frac{1}{z-w}\]

is Koszul dual to the $\beta\gamma$ system with $Q = 1$:
\[\beta(z)\gamma(w) \sim \frac{1}{z-w}\]

At positive genus the fermionic side is governed by Pfaffian line
bundles and spin parity, while the $\beta\gamma$ side is a bosonic
free-field system.  The free-field exactness result used here is the
vanishing of the mixed-channel graph sum in
Chapter~\ref{ch:genus-expansions}; no covering statement for the full
cohomology of $\overline{\mathcal M}_{g,n}$ is asserted here.

\subsubsection{Example 2: $\mathcal{W}$-algebras and their duals}

For the principal $\mathcal{W}$-algebra $\mathcal{W}^k(\mathfrak{g})$, the
chiral Koszul dual has the same modular characteristic as the
companion $\mathcal{W}^{k'}(\mathfrak{g})$, where
$k' = -k - 2h^\vee$ is the Feigin--Frenkel companion level
(Chapter~\ref{chap:w-algebra-koszul}). The stronger same-family
algebra identification is conditional on DS/bar transport. This is
chiral Koszul characteristic transport within the
$\mathcal{W}$-algebra family, distinct from the
Feigin--Frenkel--Langlands correspondence, which relates the
critical-level center to opers on $\mathfrak{g}^\vee$).
On the scalar diagonal this gives
\[
F_g^{\mathrm{diag}}(\mathcal W^k(\mathfrak g))
+F_g^{\mathrm{diag}}(\mathcal W^{k'}(\mathfrak g))
 =
\bigl(\kappa(\mathcal W^k(\mathfrak g))
+\kappa(\mathcal W^{k'}(\mathfrak g))\bigr)
\lambda_g^{\mathrm{FP}}.
\]
For principal $\mathcal W_N$ with $N\geq 3$, this is only the
diagonal scalar part; the full free energy includes the cross-channel
correction of Theorem~\ref{thm:multi-weight-genus-expansion}.

\section{Dictionary}

\begin{center}
\begin{tabular}{lll}
\toprule
Geometric term & Algebraic term & Scope \\
\midrule
Arnold relations & MC flatness & genus-$0$ collision stratum \\
Quantum corrections & stable-graph terms & positive-genus modular stratum \\
Configuration spaces & chiral bar complex & ordered points on $X$ \\
Prime form & propagator kernel & genus-$g$ curve \\
Period matrix & Hodge-bundle coordinate & scalar curvature projection \\
Koszul partner $\cA^!$ & Verdier dual of the bar coalgebra &
not $\Omega B(\cA)=\cA$ \\
\bottomrule
\end{tabular}
\end{center}

\section{The modular quantum \texorpdfstring{$L_\infty$}{L-infinity} completion}
\label{sec:modular-quantum-Linfty}
\index{modular quantum L-infinity@modular quantum $L_\infty$ completion}

The constructions of this chapter (curved $A_\infty$ operations, genus
corrections, and the quantum master equation) assemble into a single
algebraic object: the \emph{modular quantum $L_\infty$ completion} of a
chiral algebra, the cyclic/modular bar-cobar machine enriched by
the BV operator and spectral line-operator data.

\subsection{The full modular homotopy object}

\begin{definition}[Modular homotopy package]
\label{def:modular-homotopy-package}
\index{modular homotopy package}
Let $\cA$ be a chiral algebra on a curve $X$ with cyclic pairing, and
let $\mathcal L_\cA$ be its spectral line-operator algebra.  On the
chiral Koszul locus there is a Verdier comparison
$\mathcal L_\cA\simeq \cA^!$; outside that locus the symbol
$\mathcal L_\cA$ is retained.  A spectral twisting datum is an element
$r(z)\in
\mathcal L_\cA\widehat\otimes
\mathcal L_\cA((z^{-1}))$
(Definition~\ref{def:spectral-quantum-group}). The \emph{modular
homotopy package} is
\[
 \mathcal{H}(T, B_\infty)
 \;:=\;
 \bigl(\Omega^{\mathrm{mod}}\, F^{\mathrm{cyc}}\, \bar{B}(\cA_\partial),\;
 D\bigr),
\]
where $D = d_{A_\infty} + d_{\mathrm{sew}} + \hbar\,\Delta + d_r$ is
the total differential, with the four summands playing distinct
geometric roles:
\begin{enumerate}
\item $d_{A_\infty}$ is the bar differential encoding the $A_\infty$
 operations $\{m_n\}_{n \ge 0}$;
\item $d_{\mathrm{sew}}$ is the sewing differential that glues output
 legs to input legs using the cyclic pairing and the homotopy
 propagator~$h$;
\item $\Delta$ is the BV loop-contraction operator, raising the genus
 by one;
\item $d_r$ inserts the spectral line-operator twisting datum
 $r(z)$.
\end{enumerate}
This is the cyclic/modular bar-cobar machine with spectral twisting;
the four summands are kept separate because they have different
normalizations and different genus behavior.
\end{definition}

\begin{remark}[Coderivation model of the modular $L_\infty$ structure]
\label{rem:coderivation-linfty-model}
The four-component differential
$D = d_{A_\infty} + d_{\mathrm{sew}} + \hbar\Delta + d_r$
is the \emph{coderivation model} of the modular $L_\infty$ structure
on~$\mathfrak{g}_{\cA}^{\mathrm{mod}}$. Its components are the $L_\infty$ structure maps
evaluated on specific graph types: $d_{A_\infty}$ encodes tree graphs
(genus~$0$), $d_{\mathrm{sew}}$ encodes self-gluing (separating nodes),
$\hbar\Delta$ encodes non-separating clutching (genus-raising loops),
and $d_r$ encodes spectral twisting (line-operator insertions).

The bar-cobar adjunction of Theorem~A is compatible with this
coderivation model through the completed convolution formalism.
Homotopy transfer (Appendix~\ref{app:homotopy-transfer}) transports
the tree operations, while the positive-genus terms remain in the
stable-graph completion.
\end{remark}

\subsection{Tree-level operations: explicit \texorpdfstring{$A_\infty$}{A-infinity} and \texorpdfstring{$L_\infty$}{L-infinity} brackets}

At genus zero the modular package reduces to tree-level homotopy
transfer. Let $(V, d_V, i, p, h)$ be a contraction of the BV complex
onto its cohomology.

\begin{definition}[Tree-level $A_\infty$ operations]
\label{def:tree-level-ainfty}
\index{A-infinity operations@$A_\infty$ operations!tree-level transfer}
For each $n \ge 2$, the transferred $A_\infty$ operation is
\begin{equation}\label{eq:tree-mn}
 m_n^{\mathrm{tr}}
 \;=\;
 \sum_{T \in \mathrm{PBT}_n} (\pm)\, p \circ \mu_T \circ i^{\otimes n},
\end{equation}
where $\mathrm{PBT}_n$ denotes the set of planar binary rooted trees
with $n$ leaves. Internal edges carry the homotopy~$h$, each vertex
carries the interaction, each leaf carries~$i$, and the root carries~$p$.
The signs are determined by the Koszul rule with respect to the
cohomological grading.
\end{definition}

\begin{definition}[Tree-level $L_\infty$ brackets]
\label{def:tree-level-linfty}
\index{L-infinity brackets@$L_\infty$ brackets!antisymmetrization from $A_\infty$}
Antisymmetrizing yields the classical $L_\infty$ brackets:
\begin{equation}\label{eq:tree-ell}
 \ell_n^{\mathrm{cl}}(x_1, \dotsc, x_n)
 \;=\;
 \sum_{\sigma \in S_n} \chi(\sigma)\,
 m_n^{\mathrm{tr}}(x_{\sigma(1)}, \dotsc, x_{\sigma(n)}),
\end{equation}
where $\chi(\sigma)$ is the Koszul sign of the permutation.
\end{definition}

\begin{remark}[Two-level convention for quantum corrections]
\label{rem:quantum-corrections-two-level}
In the two-level convention
(\S\ref{subsec:two-level-convention}), the tree-level
$L_\infty$-brackets $\ell_n^{\mathrm{cl}}$ are $\Convinf$ at
genus~$0$; the genus-refined operations $\ell_{g,n}$ of
Definition~\ref{def:genus-refined-linfty} below are the
genus-truncation of $\Definfmod(\cA)$.
\end{remark}

\subsection{Quantum/modular completion: wheeled graphs}

Beyond genus zero, the completed operations are indexed by stable
graphs.

\begin{definition}[Genus-refined $L_\infty$ operations; \ClaimStatusProvedHere]
\label{def:genus-refined-linfty}
\index{L-infinity brackets@$L_\infty$ brackets!genus-refined}
Define the genus-$g$, $n$-ary operation by the completed stable-graph
sum
\begin{equation}\label{eq:genus-ell}
 \ell_{g,n}(x_1, \dotsc, x_n)
 \;=\;
 \sum_{\Gamma \in \mathcal{G}_{g,n}^{\mathrm{st}}}
 \frac{\hbar^g}{|\operatorname{Aut}\Gamma|}\;
 W_\Gamma\; \mathrm{Cont}_\Gamma(x_1, \dotsc, x_n),
\end{equation}
where:
\begin{itemize}
\item classical vertices carry the tree-level operations
 $m_k^{\mathrm{tr}}$ or $\ell_k^{\mathrm{cl}}$;
\item each internal edge carries the propagator~$h$;
\item internal loops are contracted via the BV pairing;
\item $W_\Gamma$ is the geometric Feynman weight, and
 $\mathrm{Cont}_\Gamma$ is the algebraic contraction.  Assign to each
 internal edge~$e$ the propagator
 $h_e = K^{(g_e)}(z_{s(e)},z_{t(e)})$ and to each vertex~$v$ the
 fundamental chain or Hodge insertion~$\omega_v$ dictated by its
 stable-curve label.  In the unintegrated tree-transfer case
 $\omega_v=1$.  Then
 \[
 W_\Gamma
 \;:=\;
 \prod_{e \in E_{\mathrm{int}}(\Gamma)} h_e
 \cdot
 \prod_{v \in V(\Gamma)} \omega_v,
 \]
 while $\mathrm{Cont}_\Gamma$ inserts the vertex operations
 $m_{|v|}^{\mathrm{cl}}$ or $\ell_{|v|}^{\mathrm{cl}}$ and pairs the
 half-edges through the cyclic/BV pairing.
	 When $\Gamma$ has internal loops ($g(\Gamma) > 0$),
	 the loop contractions $\sum_a e_a \otimes e^a$ produce
	 the BV operator~$\Delta$.  The expression is formal in the
	 completed genus/weight filtration.  Under the local-finiteness
	 hypotheses of \S\ref{subsec:hs-sewing}, each fixed
	 \((g,n)\)-weight component has finitely many contributing pole
	 contractions.  The all-genus graph series is used only as a
	 formal completed object here.
\end{itemize}
\end{definition}

\subsection{The quantum master equation}

\begin{theorem}[Quantum $L_\infty$ master equation; \ClaimStatusProvedHere]
\label{thm:quantum-linfty-master}
\index{quantum master equation!modular $L_\infty$}
The genus-refined operations $\{\ell_{g,n}\}$ satisfy the modular
$L_\infty$ constraint: for all $g \ge 0$ and $n \ge 0$ with
$2g - 2 + n > 0$,
\begin{equation}\label{eq:quantum-master-linfty}
 \frac{1}{2}
 \sum_{\substack{g_1 + g_2 = g \\ I \sqcup J = [n]}}
 (\pm)\; \ell_{g_1,\,|I|+1}\bigl(
 \ell_{g_2,\,|J|+1}(x_J),\, x_I\bigr)
 \;+\;
 \Delta\, \ell_{g-1,\,n+2}(x_1, \dotsc, x_n)
 \;=\; 0.
\end{equation}
The BV term is omitted when $g=0$.
\end{theorem}

\begin{proof}
The first sum is the quadratic genus-splitting term; the factor
$\tfrac{1}{2}$ is the Maurer--Cartan/QME factor from
Appendix~\ref{app:signs}.  The second term is the
loop-contraction/BV term.  Together they are the
$(g,n)$-component of the square-zero condition for the full modular
coderivation~$D$
(Definition~\ref{def:modular-homotopy-package}).
The identity~$D^2 = 0$ is
Theorem~\ref{thm:convolution-d-squared-zero}, which
follows from the boundary-of-boundary relation~$\partial^2 = 0$
on the stratification of~$\overline{\mathcal{M}}_{g,n}$.
The $(g,n)$-component of~$D^2 = 0$ is
precisely~\eqref{eq:quantum-master-linfty}: the
genus-splitting sum encodes self-gluing along separating
nodes and the BV term encodes gluing along non-separating
nodes.
\end{proof}

\begin{remark}[Feynman transform and BV comparison]
\label{rem:qc-feynman-bv-comparison}
The graph sum in Definition~\ref{def:genus-refined-linfty} is a
Getzler--Kapranov Feynman-transform sum for
$\mathsf F\mathrm{Com}$, rather than an analytic perturbative
integral.
The quadratic term in~\eqref{eq:quantum-master-linfty} is the
separating clutching boundary
$\overline{\mathcal M}_{g_1,|I|+1}\times
\overline{\mathcal M}_{g_2,|J|+1}\to
\overline{\mathcal M}_{g,n}$, and the BV term is the
non-separating clutching boundary
$\overline{\mathcal M}_{g-1,n+2}\to\overline{\mathcal M}_{g,n}$.
The notation $W_\Gamma$ records the same stable-graph combinatorics as
the Feynman transform.  Analytic convergence of an all-loop path
integral and comparison with a physical determinant require separate
input.
\end{remark}

\begin{remark}[Relation to the full modular homotopy type]
\label{rem:modular-homotopy-type}
\index{modular homotopy type!spectral enrichment}
The full modular homotopy type combines four layers:
\begin{enumerate}
\item the \emph{cyclic/modular bar-cobar complex}, encoding
 stable-curve gluing via the Feynman transform;
\item \emph{Verdier--Koszul duality}, when $\cA$ lies on the chiral
 Koszul locus, producing the partner $\cA^!$ as the Verdier dual of
 the bar coalgebra; the bar-cobar equivalence
 $\Omega^{\mathrm{ch}}B^{\mathrm{ch}}(\cA)\simeq \cA$ is the
 inversion theorem and is kept separate;
\item the \emph{derived open/closed center}
 $Z_{\mathrm{der}}(\cA_\partial)$, the mathematical closed-sector object
 computed by chiral Hochschild cochains of the boundary under
 Theorem~H; a physical closed-string state space requires an
 additional comparison theorem;
\item the \emph{spectral $r(z)$-defect datum}, encoding the OPE of
 line operators through $\mathcal L_\cA$; on the comparison surfaces
 where Definition~\ref{def:dg-shifted-yangian} applies, this is the
 dg-shifted-Yangian datum.
\end{enumerate}
The modular envelope sees stable-curve gluing; the addition of $r(z)$
records the spectral OPE of line operators.  On the
Dimofte--Niu--Py comparison surface this is the dg-shifted-Yangian
package of Definition~\ref{def:dg-shifted-yangian}.  Its consistency
condition is the square-zero equation for the $d_r$-twisted modular
bar complex; the Yang--Baxter extension remains the content of
Conjecture~\ref{conj:modular-yang-baxter}.
\end{remark}

\section{Non-renormalization and the exactness of tree-level operations}
\label{sec:non-renormalization-tree}
\index{non-renormalization!tree-level exactness}

Under the hypotheses of the non-renormalization theorem, the modular
homotopy package splits
into a tree-level piece (the $A_\infty$ operations $\{m_n\}$ and the
spectral kernel $r(z)$) and a loop piece (the BV operator $\Delta$
and the genus-raising operations $\ell_{g,n}$).

\begin{definition}[Quasi-linear HT theory]
\label{def:quasi-linear}
\index{quasi-linear theory}
A holomorphic-topological QFT is \emph{quasi-linear} if the
interaction vertices are at most linear in the fields being
propagated; that is, each vertex has at most one propagating leg.
The HT-twisted $3$d $\cN=2$ examples used below with polynomial
superpotential are quasi-linear when the cubic vertex
$\operatorname{Tr}(c\,Z_1 Z_2)$ is linear in the gauge field~$c$
being propagated, and the matter couplings satisfy the same
one-propagating-leg condition.
\end{definition}

\begin{theorem}[Non-renormalization at tree level;
 \ClaimStatusProvedElsewhere]
\label{thm:non-renormalization-tree}
\index{non-renormalization!Dimofte--Niu--Py}
\textup{(Dimofte--Niu--Py \cite{DNP2025}, Theorem~4.1 and
Proposition~6.1.)}
Let $\cA$ be a quasi-linear $3$d HT QFT. Then:
\begin{enumerate}
\item[\textup{(i)}] The $A_\infty$ operations $\{m_n\}$ on both
 $\cA$ and $\cA^!$ are determined at tree level: the loop diagrams
 covered by the theorem have zero contribution to $m_n$.
\item[\textup{(ii)}] The line-operator OPE kernel $r(z)$ is
 determined by tree-level ladder diagrams.
\item[\textup{(iii)}] $r(z)$ is independent of the bulk field
 interactions; it depends only on the propagator and the
 boundary coupling.
\end{enumerate}
\end{theorem}

\begin{remark}[Validation of tree-level computations]
\label{rem:non-renorm-validation}
\index{non-renormalization!validation of homotopy transfer}
For theories satisfying Theorem~\ref{thm:non-renormalization-tree},
homotopy transfer computes the tree operations
(Definition~\ref{def:tree-level-ainfty}) without additional loop
vertices.  The positive-genus terms in this chapter enter through the
genus-raising BV operator~$\Delta$ and the Hodge curvature, while the
local tree vertices stay fixed.
\end{remark}

\begin{corollary}[Collision residue normalization for standard-family
$r$-matrices;
 \ClaimStatusProvedElsewhere]
\label{cor:exact-r-matrix}
\index{r-matrix@$r$-matrix!exactness for standard families}
For the standard families $\cA \in \{\mathcal{H}_k,\;
\beta\gamma/bc,\; \widehat{\mathfrak{g}}_k \text{ at non-critical
level},\; \mathrm{Vir}_c,\; \mathcal{W}_N\}$, the $r$-matrices
computed in Part~\ref{part:characteristic-datum} via the Laplace
transform of the $\lambda$-bracket are local genus-$0$ collision
residues.  This is a statement about the singular OPE coefficient,
separate from all-genus analytic non-renormalization statements for
standard VOAs.  For the quasi-linear HT comparison theories,
Theorem~\ref{thm:non-renormalization-tree} explains why no additional
loop vertex changes this local residue.
For affine Kac--Moody algebras this means
$r^{\mathrm{coll}}(z)=k\Omega_{\mathrm{tr}}/z$, rather than the
KZ-normalized coefficient $\Omega_{\mathrm{KZ}}/((k+h^\vee)z)$.
\end{corollary}

\begin{remark}[Trees and loops]
\label{rem:trees-loops-separation}
\index{non-renormalization!trees vs.\ loops}
For the quasi-linear theories covered by
Theorem~\ref{thm:non-renormalization-tree}, the genus tower (the
$\ell_{g,n}$ operations of Definition~\ref{def:genus-refined-linfty})
is separate from the local tree-vertex data:
\[
 \text{trees} = \text{chirality},
 \qquad
 \text{loops} = \text{modularity}
\]
under exactly those hypotheses.
\end{remark}

\begin{remark}[Consequences for the modular package]
\label{rem:non-renorm-modular-consequences}
\index{non-renormalization!modular package consequences}
Under the hypotheses of Theorem~\ref{thm:non-renormalization-tree},
the modular package separates as follows:
\begin{enumerate}[label=\textup{(\roman*)},nosep]
\item \emph{The PVA $\lambda$-bracket is the local tree datum.} For
 the standard-family comparison surfaces listed above, the classical
 $\lambda$-bracket
 (Definition~\ref{def:affine-pva-lambda-bracket} for affine,
 \S\ref{subsec:virasoro-teichmuller-phase-space} for Virasoro,
 \S\ref{sec:w3-pva-holography} for $\mathcal{W}_3$) is the
 tree-level datum used by the local reconstruction of the collision
 residue.
\item \emph{The genus-$0$ bar differential is unchanged.} For
 quasi-linear theories,
 the bar differential $d_{\bar{B}}$ at genus zero receives no
 loop-vertex correction.  The genus-raising terms enter through the
 curvature $m_{0,\mathrm{scal}}^{(g)}=\kappa(\cA)\lambda_g$ on the
 scalar lane and through cross-channel graph sums for interacting
 multi-weight algebras.
\item \emph{The collision $r$-matrix remains a tree datum.} The
 line-operator OPE kernel $r(z)$ computed in
 Part~\ref{part:characteristic-datum} is the collision residue
 normalization; it is distinct from the KZ coefficient in
 \eqref{eq:qc-trace-kz-normalizations}.
\end{enumerate}
The genus filtration of the master MC element has the decomposition
\[
 \Theta_{\cA}
 \;=\;
 \underbrace{\Theta_{\cA}\big|_{\hbar=0}}_{\text{tree: }\tau}
 \;+\;
 \underbrace{\Theta_{\cA}\big|_{\hbar^g,\;g \geq 1}}_{\text{loop: genus tower}},
\]
with the tree part $\tau$ the twisting morphism of Theorem~A.  This is
a filtration statement for the completed MC element.  Analytic
summability of the all-genus expansion and reconstruction of the full
partition function require separate input.
\end{remark}

\section{The Jacobiator homotopy and the Eilenberg--Zilber operad}
\label{sec:jacobiator-homotopy}
\index{Jacobiator!homotopy}
\index{Eilenberg--Zilber operad}

The homotopy chiral algebra structure
(Theorem~\ref{thm:cech-hca}) equips the \v{C}ech resolution of a
chiral algebra with higher operations $\{F_n\}_{n \geq 1}$
satisfying an $L_\infty$ master equation. At degree~$3$, this
master equation has a transparent interpretation: it is the
\emph{Jacobiator homotopy}, the precise mechanism by which Jacobi
fails on the resolution and is restored by a coherent correction.
The Jacobiator appears only for covers with sufficient overlap: on
$\mathbb{P}^1$ with the standard two-element cover, it vanishes
identically. The obstruction requires triple intersections (three
or more overlapping opens) and thus measures genuine descent failure
for the chosen cover.

The Jacobiator, and all higher
$F_n$, are components of a single master equation
$D \cdot \Phi + \tfrac{1}{2}[\Phi,\Phi] = 0$, where
$\Phi = \sum_{n \geq 2} F_n$ is the MC element encoding the
entire HCA structure. By Proposition~\ref{prop:chriss-ginzburg-structure}, this is a
projection of the master equation
$D_{\cA}\Theta_{\cA} + \tfrac{1}{2}[\Theta_{\cA},\Theta_{\cA}] = 0$
to the \v{C}ech site at genus~$0$.

\subsection{The Jacobiator on the \v{C}ech complex}
\label{subsec:jacobiator-cech}

Let $\cA$ be a chiral algebra on~$X$ with cover $\mathcal{U}$,
and let $[{-},{-}] = F_2$ denote the binary chiral bracket on the
Moore complex $M\check{C}^\bullet(\mathcal{U};\cA)$. The Jacobi
identity holds \emph{strictly} on each open $U_{i_0\cdots i_p}$
(where the full chiral product is defined), but on the total
\v{C}ech complex the bracket interacts with the \v{C}ech
differential, and the Jacobi identity may fail.

\begin{definition}[Jacobiator]
\label{def:jacobiator}
For $a,b,c \in M\check{C}^\bullet$, the \emph{Jacobiator} is
\[
\mathrm{Jac}(a,b,c)
= [a,[b,c]] + (-1)^{|a|(|b|+|c|)}[b,[c,a]]
 + (-1)^{|c|(|a|+|b|)}[c,[a,b]].
\]
On the total \v{C}ech complex, $\mathrm{Jac}$ is in general
\emph{nonzero}: it measures the obstruction to $F_2$ being a
strict Lie bracket on the resolution.
\end{definition}

\begin{proposition}[Two-element covers are strict;
\ClaimStatusProvedHere]
\label{prop:two-element-strict}
\index{Jacobiator!vanishing for two-element covers}
Let $\mathcal{U} = \{U_0, U_\infty\}$ be a two-element affine
open cover of~$X$ (e.g., the standard cover of
$\mathbb{P}^1$). Then
$\check{C}^p(\mathcal{U};\cA) = 0$ for $p \geq 2$,
and consequently:
\begin{enumerate}[label=\textup{(\roman*)}]
\item The Jacobiator vanishes:
 $\mathrm{Jac}(a,b,c) = 0$ for all $a,b,c$.
\item All higher operations vanish: $F_n = 0$ for $n \geq 3$.
\item The \v{C}ech HCA is \emph{strict}: the binary bracket
 $F_2$ is a genuine Lie bracket on $M\check{C}^\bullet$.
\end{enumerate}
\end{proposition}
\begin{proof}
A two-element cover has $\check{C}^0 = \cA(U_0) \oplus
\cA(U_\infty)$, $\check{C}^1 = \cA(U_0 \cap U_\infty)$,
and $\check{C}^p = 0$ for $p \geq 2$. In particular, the
Moore complex is concentrated in degrees~$0$ and~$1$. The
$n$-ary operation $F_n$ raises total \v{C}ech degree by $n-2$
(it maps $n$ inputs summing to degree~$m_1 + \cdots + m_n$
to degree $m_1 + \cdots + m_n$, while the chain in
$\mathcal{Z}(n)$ contributes degree $2-n$; see
Remark~\ref{rem:operadic-action-map}). For $n \geq 3$, the
output lands in $\check{C}^p$ for $p \geq 2$, which vanishes.
Thus $F_n = 0$ for all $n \geq 3$. In particular the Jacobiator,
which is the boundary of~$F_3$
(equation~\eqref{eq:jacobiator-homotopy}), vanishes.
\end{proof}

\begin{remark}[When the Jacobiator appears]
\label{rem:jacobiator-when-nonzero}
\index{Jacobiator!triple intersection origin}
The vanishing of Proposition~\ref{prop:two-element-strict} is
\emph{not} a feature of the algebra: it is a feature of the
\emph{cover}. For the \emph{same} chiral algebra on the same
curve~$X$ but with a cover $\mathcal{U}' = \{U_0,U_1,U_2\}$
having nontrivial triple intersection
$U_0 \cap U_1 \cap U_2 \neq \varnothing$, the \v{C}ech complex
extends to degree~$2$ and the Jacobiator can be nonzero. The
correction $F_3$ then measures the \emph{descent obstruction}:
the chiral bracket holds strictly on each open, but gluing
these local brackets across a triple overlap requires a
coherent homotopy.

For $\widehat{\mathfrak{sl}}_2$ at level~$k$ on $\mathbb{P}^1$,
the two-element cover gives a strict HCA
(Example~\ref{ex:cech-hca-heisenberg}); but refining to a
three-element cover (e.g., the standard affine cover of a
projective curve of genus $g \geq 1$) produces a nonzero
Jacobiator from the non-abelian OPE $[e,f] = h$
(Example~\ref{ex:cech-hca-sl2}). This is the prototypical
example: the Jacobiator is a descent obstruction for the cover, while
the local OPE supplies its Lie input.
\end{remark}

\subsection{The Jacobiator as the degree-$3$ master equation}
\label{subsec:jacobiator-master}

The Jacobiator homotopy is the degree-$3$ component of the MC equation
$D \cdot \Phi + \tfrac{1}{2}[\Phi,\Phi] = 0$
governing the HCA structure map
$\Phi = \sum_{n \geq 2} F_n$. The degree-$n$ component reads:
\[
d_M \circ F_n
+ F_n \circ d_M^{\otimes n}
+ \sum_{\substack{p+q=n+1 \\ p,q \geq 2}}
 F_p \circ_{(i)} F_q
= 0,
\]
where $\circ_{(i)}$ denotes operadic composition at the $i$-th
slot, summed over all $i$ and all $(p,q)$-shuffles. At
degree~$3$, the quadratic term reduces to $F_2 \circ F_2$
(iterated binary bracket), recovering:

\begin{proposition}[Jacobiator is nullhomotopic;
\ClaimStatusProvedElsewhere]
\label{prop:jacobiator-nullhomotopic}
\index{Jacobiator!nullhomotopy}
\textup{(Malikov--Schechtman \cite{MS24}, \S4.1.)}
There exists a trilinear operation
$J'\colon (M\check{C})^{\otimes 3} \to M\check{C}[-1]$ such that
\begin{equation}\label{eq:jacobiator-homotopy}
d_M \circ J' + J' \circ d_M^{\otimes 3} = \mathrm{Jac},
\end{equation}
where $d_M^{\otimes 3}$ is the Leibniz extension of $d_M$ to the
triple tensor product. More precisely, $J'$ is the ternary component
$F_3$ of the HCA structure
(Theorem~\textup{\ref{thm:cech-hca}}), and
equation~\eqref{eq:jacobiator-homotopy} is the $L_\infty$
relation at degree~$3$: the degree-$3$ projection of the master
equation $D \cdot \Phi + \tfrac{1}{2}[\Phi,\Phi] = 0$.
\end{proposition}

\subsection{Translation to bar-side $A_\infty$ operations}
\label{subsec:jacobiator-bar}

The bar-side operation $m_3$ and the \v{C}ech-side homotopy
$F_3$ encode the same degree-$3$ correction, but through different
geometric mechanisms. On the bar side, $m_3$ arises from
integration over the Fulton--MacPherson compactification
$\overline{\mathrm{FM}}_3(\mathbb{C})$: it is the homotopy-transfer
datum obtained by composing binary interactions through the
propagator on the configuration space of three points. On the
\v{C}ech side, $F_3$ arises from triple intersections of the cover:
it measures how local Jacobi identities fail to descend.

The \v{C}ech--bar comparison
$\Phi\colon (M\check{C},\{F_n\}) \xrightarrow{\sim}
(\barBgeom(\cA),\{m_n\})$
(Conjecture~\ref{conj:cech-bar-intertwining}) translates between
these two pictures. Under~$\Phi$:
\begin{itemize}
\item $\Phi$ intertwines $F_3$ with $m_3$: integration over
 $\overline{\mathrm{FM}}_3(\mathbb{C})$ on the bar side
 corresponds to the triple-intersection descent correction on the
 \v{C}ech side.
\item The $A_\infty$ relation at degree~$3$,
\[
m_1 \circ m_3 + m_3 \circ (m_1 \otimes 1 \otimes 1
 + 1 \otimes m_1 \otimes 1 + 1 \otimes 1 \otimes m_1)
+ m_2 \circ (m_2 \otimes 1 + 1 \otimes m_2) = 0,
\]
is the bar-side avatar of~\eqref{eq:jacobiator-homotopy}:
the first two terms correspond to
$d_M \circ F_3 + F_3 \circ d_M$,
while $m_2(m_2 \otimes 1 + 1 \otimes m_2)$ is the Jacobiator
$\mathrm{Jac}$ expressed through iterated binary products.
\item More generally, $\Phi$ intertwines the full MC equations on
 both sides: the \v{C}ech master equation
 $D \cdot \Phi_{\check{C}} + \tfrac{1}{2}
 [\Phi_{\check{C}},\Phi_{\check{C}}] = 0$ maps to
 the bar master equation encoded by $\{m_n\}$.
\end{itemize}

\subsection{The Eilenberg--Zilber operad as genus-$0$ shadow of
$\mathsf{F}\mathrm{Com}$}
\label{subsec:ez-fcom}
\index{Eilenberg--Zilber operad!genus-0 shadow of FCom@genus-$0$
shadow of $\mathsf{F}\mathrm{Com}$}

\begin{remark}[Operadic control at genus~$0$ and beyond]
\label{rem:ez-operad-controls-ainfty}
The Eilenberg--Zilber operad
$\mathcal{Y}(n) = \mathcal{Z}(n) \otimes_k \mathrm{Lie}(n)$
of Malikov--Schechtman~\cite{MS24} and the Feynman transform
$\mathsf{F}\mathrm{Com}$ of the commutative modular operad
(Theorem~\ref{thm:bar-modular-operad}) control homotopies at
different levels:
\begin{center}
\renewcommand{\arraystretch}{1.2}
\begin{tabular}{lll}
\toprule
& \emph{\v{C}ech side} & \emph{Bar side} \\
\midrule
Genus $0$ & Eilenberg--Zilber $\mathcal{Y}(n)$ &
 Associahedra $K_n \subset \overline{\mathrm{FM}}_n(\mathbb{C})$ \\
Genus $g \geq 1$ & Cosimplicial modular operad &
 $\mathsf{F}\mathrm{Com}$, stable graphs on
 $\overline{\mathcal{M}}_{g,n}$ \\
\bottomrule
\end{tabular}
\end{center}
At genus~$0$, the chain complex
$\mathcal{Z}(n) = \operatorname{Tot}\operatorname{Hom}_\Delta
(Z, Z^{\otimes n})$ encodes the combinatorics of the simplicial
structure (how sections on intersections compose), while
$\mathrm{Lie}(n)$ encodes the chiral bracket. Their tensor product
$\mathcal{Y}(n)$ maps quasi-isomorphically to $\mathrm{Lie}(n)$
via the augmentation
$\varepsilon\colon \mathcal{Z}(n) \to k$,
so the HCA is a resolution of a strict chiral algebra.

At genus $g \geq 1$, the \v{C}ech homotopies no longer suffice:
the bar complex sees the modular operad, and the homotopies are
controlled by stable-graph compositions in
$\mathsf{F}\mathrm{Com}$.  The Eilenberg--Zilber operad
$\mathcal{Y}(n)$ is the simplicial tree model feeding the genus-$0$
part of this machine.  It records parenthesized collision and descent
data; the full genus-$0$ summand of $\mathsf{F}\mathrm{Com}$ also
records the complex cross-ratio moduli of pointed rational curves.
The passage from
$\mathcal{Y}(n)$ to $\mathsf{F}\mathrm{Com}$ is therefore a
filtered comparison from tree-level \v{C}ech homotopies to the
stable-graph bar homotopies, paralleling the genus expansion
$\Theta_{\cA} = \Theta_{\cA}|_{\hbar=0} +
\sum_{g \geq 1}\hbar^g\,\Theta_{\cA,g}$ of the master MC element.
\end{remark}

\begin{remark}[Tree comparison: EZ chain complex, associahedra,
and genus-$0$ $\mathsf{F}\mathrm{Com}$]
\label{rem:ez-fcom-precise-truncation}
\index{Eilenberg--Zilber operad!precise truncation to FCom@precise
truncation to $\mathsf{F}\mathrm{Com}$}
\index{associahedron!and genus-0 FCom@and genus-$0$
$\mathsf{F}\mathrm{Com}$}
The comparison between $\mathcal{Y}(n)$ and
$\mathsf{F}\mathrm{Com}$ is a filtered tree comparison.  A single
associahedron sees the tree chamber, while
$\overline{\mathcal M}_{0,n+1}$ carries the full pointed-rational
moduli complex.

\smallskip\noindent\emph{(i) Genus-$0$ component of
$\mathsf{F}\mathrm{Com}$.}
The space $\Conf_n(\mathbb{C})/\mathrm{Aff}$, where
$\mathrm{Aff} = \mathbb{C} \rtimes \mathbb{C}^\times$ acts by
affine transformations, is the open genus-$0$ configuration chart.
Its stable compactification is $\overline{\mathcal{M}}_{0,n+1}$.
The cellular chains
$C_*(\overline{\mathcal{M}}_{0,n+1};\, k)$ form the genus-$0$
component of the Feynman transform:
\[
\mathsf{F}\mathrm{Com}(0,\, n{+}1)
= C_*(\overline{\mathcal{M}}_{0,n+1};\, k),
\]
with operadic composition induced by the clutching maps
$\overline{\mathcal{M}}_{0,n_1+1} \times
\overline{\mathcal{M}}_{0,n_2+1}
\to \overline{\mathcal{M}}_{0,n_1+n_2}$
along boundary divisors.

\smallskip\noindent\emph{(ii) EZ chain complex and the
associahedron.}
The Eilenberg--Zilber chain complex
$\mathcal{Z}(n)
= \operatorname{Tot}\operatorname{Hom}_\Delta(Z,\, Z^{\otimes n})$
carries a natural filtration by simplicial degree. The Stasheff
associahedron $K_n$ (the convex polytope whose face poset encodes
all ways of parenthesizing $n$ letters) provides a geometric
model: its cellular chains $C_*(K_n;\, k)$ compute the same
$A_\infty$ homotopies. The comparison is the classical
quasi-isomorphism of Eilenberg--Zilber--Loday type:
\[
\alpha_n \colon \mathcal{Z}(n)
 \xrightarrow{\;\sim\;} C_*(K_n;\, k),
\]
which sends the total complex of cosimplicial Hom to the cellular
chains on the associahedron by taking the acyclic-models
comparison between the bar construction of the simplicial
structure and the CW decomposition of $K_n$.

\smallskip\noindent\emph{(iii) Associahedra inside
$\overline{\mathcal{M}}_{0,n+1}$.}
The Stasheff associahedron $K_n$ occurs in the real locus and in the
boundary-stratified collision model: dihedral orderings of the marked
points give associahedral chambers.  Choosing one chamber gives a
chain map
\[
\beta_n \colon C_*(K_n;\, k)
 \longrightarrow C_*(\overline{\mathcal{M}}_{0,n+1};\, k)
 = \mathsf{F}\mathrm{Com}(0,\, n{+}1).
\]
For $n+1\geq 5$, $\overline{\mathcal{M}}_{0,n+1}$ has complex moduli
and boundary cohomology beyond a single associahedron.  The map
identifies the tree-cell contribution used by homotopy transfer.

\smallskip\noindent\emph{(iv) The truncation chain.}
Composing the comparison maps gives the tree-level route:
\[
\mathcal{Z}(n)
 \xrightarrow[\text{EZ--Loday}]{\;\alpha_n\;}
C_*(K_n;\, k)
 \xrightarrow[\text{tree chamber}]{\;\beta_n\;}
\mathsf{F}\mathrm{Com}(0,\, n{+}1)
 \hookrightarrow
\mathsf{F}\mathrm{Com}(\bullet,\, n{+}1).
\]
The first arrow is the EZ--Loday comparison.  The second is the
tree-cell inclusion into the genus-$0$ moduli complex.  The last is
the inclusion of the genus-$0$ summand into the full modular operad.
Tensoring with the Lie operad recovers the Eilenberg--Zilber
tree contribution:
\[
\mathcal{Y}(n) = \mathcal{Z}(n) \otimes_k \mathrm{Lie}(n)
\xrightarrow[\sim]{\;\alpha_n \otimes \mathrm{id}\;}
C_*(K_n;\, k) \otimes_k \mathrm{Lie}(n)
\xrightarrow{\;\beta_n \otimes \mathrm{id}\;}
\mathsf{F}\mathrm{Com}(0,\, n{+}1) \otimes_k \mathrm{Lie}(n).
\]
Thus $\mathcal{Y}(n)$ sees the tree-level genus-$0$ collision
sector of $\mathsf{F}\mathrm{Com}$, tensored with the bracket
symmetry.
At genus $g \geq 1$, the summands
$\mathsf{F}\mathrm{Com}(g,\, n{+}1)$ involve stable graphs with
loops (corresponding to $\overline{\mathcal{M}}_{g,n+1}$), which
have no counterpart in the tree-level simplicial structure
underlying $\mathcal{Z}(n)$. This is the precise sense in which
the EZ operad is the \v{C}ech tree shadow of the modular
bar construction.
\end{remark}

\subsection{Secondary Borcherds operations}
\label{subsec:secondary-borcherds}
\index{Borcherds identity!secondary operations}
\index{secondary Borcherds operations|textbf}

The HCA structure of Theorem~\ref{thm:cech-hca} produces not only
the Jacobiator homotopy $J' = F_3$ but an entire tower of
\emph{secondary Borcherds operations}
(Malikov--Schechtman~\cite{MS24}, \S4.5):

\begin{definition}[Secondary Borcherds operations]
\label{def:secondary-borcherds}
For a cosimplicial translationally invariant chiral algebra
$\mathcal{V}^\bullet$ on $C = \operatorname{Spec} k[t]$, the
homotopy chiral algebra structure produces ternary
operations of degree~$-1$:
\[
j'_{(p,q,r)}\colon
 (MV^\bullet)^{\otimes 3} \longrightarrow MV^\bullet[-1],
 \qquad p,q,r \in \mathbb{Z},
\]
satisfying the \emph{secondary Borcherds identities}
\textup{(\cite{MS24}, equation~(4.5.1))}:
\begin{multline}\label{eq:secondary-borcherds}
 (d \circ_{(p,q,r)} +\; {}_{(p,q,r)} \circ d)(a \otimes b \otimes c)
 \;=\\
 \sum_{j \geq 0}\binom{p}{j}\bigl(
 (-1)^j a_{(p+q-j)} b_{(r+j)} c
 - (-1)^{j+p+|a|\,|b|}\, b_{(r+p-j)} a_{(q+j)} c
 \bigr)\\
 - \sum_{j \geq 0}\binom{q}{j}\,
 (a_{(p+q-j)} b)_{(r+j)}\, c,
\end{multline}
where $a_{(n)}b = \phi_{[,]}((a \otimes b)(t_1-t_2)^n)$
denotes the $n$-th mode product.
\end{definition}

\begin{remark}[The $n$-th product relations as descent data]
\label{rem:vertex-lie-jacobi-family}
\index{vertex Lie algebra!Jacobi family}
For translationally invariant chiral algebras on
$X = \operatorname{Spec}k[t]$, the $n$-th products
$a_{(n)}b$ satisfy the Jacobi identity family
\cite{MS24}:
\begin{equation}\label{eq:vertex-jacobi-family}
[a_{(m)}, b_{(n)}]c
= \sum_{j \geq 0} \binom{m}{j}(a_{(j)}b)_{(m+n-j)}c
\end{equation}
and the skew-symmetry relation
\begin{equation}\label{eq:vertex-skew}
a_{(n)}b = (-1)^{n+1}\sum_{j \geq 0}
\frac{(-1)^j}{j!}\,\partial^j(b_{(n+j)}a).
\end{equation}
These formulas \emph{are} the $n$-th product relations of
Beilinson--Drinfeld~\cite{BD04}; they hold strictly on each open
of the cover. On the \v{C}ech complex, however, gluing local
brackets across multiple opens can introduce a descent obstruction;
in the HCA model the homotopies correcting this obstruction are the
operations $F_n$ for $n \geq 3$. The secondary Borcherds
identity~\eqref{eq:secondary-borcherds} is the explicit formula
for this descent correction at degree~$3$: the right-hand side is
the local $n$-th product relation (which holds on each open), and
the left-hand side is the boundary of $F_3$, which records how
that relation fails to descend.

For two-element covers the descent obstruction vanishes
(Proposition~\ref{prop:two-element-strict}); the correction $F_n$
can be nonzero only if the nerve has nondegenerate
$(n{-}1)$-simplices.  Higher overlaps are the geometric support
condition; non-vanishing requires the algebra-specific contact
representative.
\end{remark}

\begin{remark}[Derivation of the vertex Lie Jacobi family from
the Borcherds identity]
\label{rem:jacobi-family-derivation}
\index{Borcherds identity!derivation of Jacobi family}
\index{vertex Lie algebra!Jacobi family derivation}
The Jacobi identity family~\eqref{eq:vertex-jacobi-family} follows
from the Borcherds
identity for the chiral product, applied to three sections on
pairwise intersections of a \v{C}ech cover. This makes precise
the passage from OPE associativity on a single open to the $n$-th
product relations on the \v{C}ech complex.

\smallskip\noindent\emph{Setup.}
Let $\mathfrak{U} = \{U_0, U_1, U_2\}$ be a cover of
$X = \operatorname{Spec}k[t]$ with pairwise intersections
$U_{ij} = U_i \cap U_j$ and triple intersection
$U_{012} = U_0 \cap U_1 \cap U_2$. Let $a,b,c$ be sections of a
translationally invariant chiral algebra $\cA$ on $X$. On each
$U_{ij}$, the chiral product
$\mu\colon j_*j^*(\cA \boxtimes \cA) \to \Delta_*\cA$
(where $j\colon X^2 \setminus \Delta \hookrightarrow X^2$)
determines the $n$-th mode products $a_{(n)}b$ via residues:
\[
a_{(n)}b = \operatorname{Res}_{z_1 = z_2}
 \mu(a(z_1) \otimes b(z_2))(z_1 - z_2)^n.
\]

\smallskip\noindent\emph{Step 1: the Borcherds identity on a
 single open.}
On $U_{012}$, OPE associativity for the chiral product gives the
\emph{Borcherds identity}: for all $m,n \in \mathbb{Z}$,
\begin{multline}\label{eq:borcherds-ope-assoc}
\sum_{j \geq 0}\binom{m}{j}
\bigl(a_{(m+n-j)}\,b_{(r+j)} - (-1)^{m+|a|\,|b|}\,
 b_{(m+r-j)}\,a_{(n+j)}\bigr)c \\
= \sum_{j \geq 0}\binom{n}{j}
 (-1)^j\,(a_{(n-j)}\,b)_{(m+r+j)}\,c.
\end{multline}
This is the triple-point constraint: it encodes that the two ways
of composing three OPEs ($(a \cdot b) \cdot c$ versus
$a \cdot (b \cdot c)$ minus the transposition) agree as
operations on $U_{012}$. The identity holds on every open of the
cover where all three sections are defined.

\smallskip\noindent\emph{Step 2: specialization to the commutator.}
Setting $r = n$ in~\eqref{eq:borcherds-ope-assoc} and extracting
the $j = 0$ term on the left-hand side (all other terms with
$j \geq 1$ contribute higher-order corrections controlled by the
binomial coefficients), one obtains the commutator formula:
\begin{equation}\label{eq:jacobi-from-borcherds}
[a_{(m)}, b_{(n)}]c
= \sum_{j \geq 0}\binom{m}{j}(a_{(j)}\,b)_{(m+n-j)}\,c,
\end{equation}
which is precisely~\eqref{eq:vertex-jacobi-family}. The key
step: on the left-hand side, the commutator
$[a_{(m)}, b_{(n)}] = a_{(m)} b_{(n)} - (-1)^{|a|\,|b|}
b_{(n)} a_{(m)}$ arises from the two orderings of evaluation
at the pairwise intersections $U_{01} \cap U_{02}$ versus
$U_{02} \cap U_{01}$; on the right-hand side, the iterated
product $(a_{(j)}b)_{(m+n-j)}c$ arises from first composing on
$U_{01}$ and then applying the result on $U_{(01)2}$.

\smallskip\noindent\emph{Step 3: the Eilenberg--Zilber correction.}
The identity~\eqref{eq:jacobi-from-borcherds} holds \emph{strictly}
on the triple intersection $U_{012}$, because the Borcherds
identity~\eqref{eq:borcherds-ope-assoc} is a local statement.
On the \v{C}ech complex, however, passing from the triple
intersection to the global Moore complex requires choosing an
ordering of the opens and applying the Eilenberg--Zilber shuffle
map
$\nabla\colon C_*(\mathfrak{U}) \otimes C_*(\mathfrak{U})
\to C_*(\mathfrak{U} \times \mathfrak{U})$.
The shuffle introduces correction terms of the form
$\nabla - \nabla_0$ (where $\nabla_0$ is the Alexander--Whitney
map), and these correction terms are precisely the secondary
Borcherds operations $j'_{(p,q,r)}$ of
Definition~\ref{def:secondary-borcherds}. Explicitly: the
boundary of $j'_{(p,q,r)}$ equals the difference between the
local Borcherds identity (which holds on each open) and the
global identity (which holds only up to homotopy on the Moore
complex), giving~\eqref{eq:secondary-borcherds}.
\end{remark}

\begin{remark}[Secondary Borcherds and $d_{\mathrm{bracket}}^2 \neq 0$]
\label{rem:secondary-borcherds-d2}
\index{d-bracket squared@$d_{\mathrm{bracket}}^2 \neq 0$!secondary Borcherds interpretation}
The proved result that $d_{\mathrm{bracket}}^2 \neq 0$ (all $2048$
signs verified; Remark~\ref{rem:frame-why-dbracket-fails})
while the full Borcherds differential gives $d^2 = 0$ is \emph{precisely}
the Ch$_\infty$ phenomenon: the bracket-level failure is the obstruction
to strict chirality, and the secondary Borcherds operations
$j'_{(p,q,r)}$ are the coherent homotopies that restore $d^2 = 0$ at
the full level. In the bar-side language
(\S\ref{subsec:jacobiator-bar}), the secondary Borcherds
identities~\eqref{eq:secondary-borcherds} are the vertex-algebra
avatar of the $A_\infty$ associativity relations for $m_3$.

The raw HCA operations depend on the cover.  After projection to the
genus-$0$ shadow quotient and on covers whose nerves support the
required intersections, the cover-independent termination profile is:
\begin{center}
\renewcommand{\arraystretch}{1.2}
\begin{tabular}{lll}
\toprule
\emph{Family} & \emph{Shadow profile} & \emph{Archetype} \\
\midrule
Heisenberg & no cubic or higher shadow class & Gaussian (strict) \\
Affine $\widehat{\mathfrak{sl}}_2$ & first nonzero class at degree $3$ & Lie/tree \\
$\beta\gamma$ & quartic contact class, free-field cross-channel vanishing & Contact/quartic \\
Virasoro/$W_N$ & no finite termination in the standard completion & Mixed (infinite tower) \\
\bottomrule
\end{tabular}
\end{center}
The table is a statement about the projected shadow tower.  Raw
\v{C}ech non-vanishing remains cover-dependent.
\end{remark}

\begin{proposition}[Secondary Borcherds operations and the first shadow
obstructions; \ClaimStatusProvedHere]
\label{prop:borcherds-shadow-identification}
\index{secondary Borcherds operations!shadow tower identification}
Under the genus-$0$ operadic equivalence between the
Malikov--Schechtman HCA operations and the transferred
$L_\infty$ brackets on the bar-side minimal model, the HCA
operation $F_3=J'$ maps to the cubic shadow obstruction class
$o_3(\cA)$
\textup{(}Definition~\textup{\ref{def:shadow-postnikov-tower})}
at genus~$0$.  The secondary Borcherds operations give the next
computable contact representatives; after projection to the quartic
shadow quotient they detect $o_4(\cA)$ when the \v{C}ech--bar
comparison is defined on that quotient.  The explicit formula
\cite[formula~(4.5.1)]{MS24} is:
\begin{multline}\label{eq:secondary-borcherds-explicit}
 \bigl(d \circ {}_{(p,q,r)} + {}_{(p,q,r)} \circ d\bigr)(a,b,c) \\
 = \sum_{j \geq 0}\tbinom{p}{j}
 \bigl((-1)^j a_{(p+q-j)}\, b_{(r+j)}\, c
 - (-1)^{j+p+|a||b|}\, b_{(r+p-j)}\, a_{(q+j)}\, c\bigr) \\
 - \sum_{j \geq 0}\tbinom{q}{j}
 \bigl(a_{(p+q-j)}\, b\bigr)_{(r+j)}\, c,
\end{multline}
where ${}_{(p,q,r)}$ is the secondary Borcherds operation at
multi-index $(p,q,r)$. This is the three-variable refinement of the
Jacobiator homotopy: the left-hand side is the differential of the
ternary secondary operation, and the right-hand side is the secondary
Jacobi identity expressed through iterated Borcherds products.

The identification is:
\begin{enumerate}[label=\textup{(\roman*)}]
\item $F_3 = J'$ \textup{(}Jacobiator homotopy\textup{)}
 corresponds to $o_3(\cA)$ \textup{(}cubic shadow obstruction\textup{)}.
\item The projected secondary Borcherds contact class corresponds to
 $o_4(\cA)$ \textup{(}quartic resonance obstruction\textup{)} when the
 genus-$0$ \v{C}ech--bar comparison is evaluated on the quartic shadow
 quotient. The Virasoro quartic contact invariant
 $\mathfrak{Q}^{\mathrm{contact}}_{\mathrm{Vir}} = 10/[c(5c+22)]$
 is the normalized coefficient of the secondary Borcherds
 contact representative evaluated on the Moore complex of
 $\mathrm{Vir}_c$.
\item At degree $n \geq 5$, the higher HCA operations supply
 candidate representatives for the higher shadow classes
 $o_n(\cA) \in H^2(\cA^{\mathrm{sh}}_n)$; equality with the
 bar-side classes requires the same genus-$0$ comparison on the
 corresponding shadow quotient.
\end{enumerate}
The identification is structural at genus~$0$; at genus $g \geq 1$,
the bar side includes modular corrections (stable graph sums with
loops) that have no \v{C}ech counterpart in~\cite{MS24}.
\end{proposition}

\begin{proof}
The proof uses only the genus-$0$ operadic equivalence.  A full
\v{C}ech--bar $L_\infty$ intertwining statement is a separate
comparison theorem.

Over a field of characteristic~$0$, the category of
$\tau_{\leq 0}\mathcal{Y}$-algebras (Malikov--Schechtman's
homotopy Lie operadic model, \cite[\S2.6]{MS24}) and the category
of $L_\infty$-algebras (the coalgebraic model,
\cite[\S10.3]{LV12}) are equivalent: a
$\tau_{\leq 0}\mathcal{Y}$-algebra structure on a complex~$V$ is
the same as an $L_\infty$-algebra structure on~$V$ with the same
MC moduli (Remark~\ref{rem:conv-functoriality-infrastructure}(iii)).
Under this equivalence, the $n$-th HCA operation $F_n$ corresponds
to the $n$-th $L_\infty$ bracket $\ell_n$ on the Moore complex.

At genus~$0$, the $L_\infty$ brackets on the bar complex are the
transferred operations from the minimal tree model of
$C_*(\overline{\cM}_{0,n+1})$
(Construction~\ref{constr:explicit-convolution-linfty}), and the
obstruction class $o_n(\cA)$ is the cohomology class
$[\ell_n] \in H^2(\cA^{\mathrm{sh}}_n)$ when the comparison map is
defined on the relevant shadow quotient.

Part~(i): $F_3 = J'$ is the Jacobiator homotopy
(Proposition~\ref{prop:jacobiator-nullhomotopic}); $\ell_3 = o_3$
is the cubic shadow (Definition~\ref{def:shadow-postnikov-tower}).
Part~(ii): the secondary Borcherds contact representative projects
to the quartic resonance class, and the Virasoro coefficient is the
canonical value recorded in the landscape census,
$10/[c(5c+22)]$.
\end{proof}

\begin{remark}[One-slot functoriality obstruction and MC3]
\label{rem:one-slot-obstruction}%
\index{one-slot obstruction}%
\index{MC3!one-slot constraint}%
The Robert-Nicoud--Wierstra theorem
(Theorem~\ref{thm:operadic-homotopy-convolution}(iii)) states that
the convolution construction $\operatorname{hom}^\alpha$ extends to
$\infty$-morphisms in either slot separately but \emph{not both simultaneously} (it is \emph{not} a bifunctor in $\infty$-morphisms).
An explicit counterexample is given in \cite[\S6]{RNW19}.

This constrains the MC3 categorical lift: extending the
prefundamental Clebsch--Gordan closure
(Proposition~\ref{prop:prefundamental-clebsch-gordan}) from the
character level (K$_0$, one slot) to the categorical level (both
slots) cannot be achieved by a single bifunctorial extension. The
categorical lift (Theorem~\ref{thm:categorical-cg-all-types}) must
therefore proceed
\emph{one slot at a time}: first extend in the coalgebra slot (via
$\operatorname{hom}^\alpha_\ell(\Phi,1)$), then in the algebra slot
(via $\operatorname{hom}^\alpha_r(1,\Psi)$), accepting that the
two extensions are compatible only up to higher homotopy.
\end{remark}

\section{Imported K3 boundary data}
\label{sec:qc-supplement}

\begin{remark}[Pentagon coboundary at $\hbar^3$]
\label{rem:qc-1}
The following formula is imported from Vol.~III as boundary data for
comparison with the K3 chiral bialgebra.  In the notation of
Chapter~\ref{V3-ch:k3-chiral-bialgebra-platonic}, the pentagon-level
correction is
\[
 \phi^{(3)}
 \;=\;
 \zeta(3)\, c_{\mathrm{symm}}
 \;+\;
 \tfrac{25}{3}\, c_{\mathrm{timelike}}
 \;+\;
 \tfrac{\Phi_{10}}{\eta^{24}}\, c_{\Phi_{10}},
\]
with three independent third-order pieces: the Drinfeld--Kohno
associator contribution $\zeta(3)c_{\mathrm{symm}}$, the timelike
coefficient $(25/3)c_{\mathrm{timelike}}$, and the Borcherds
imaginary-root contribution
$(\Phi_{10}/\eta^{24})c_{\Phi_{10}}$
\cite{Drinfeld1990,Borcherds1992,GritsenkoNikulin1998}.
\end{remark}

\begin{remark}[$\widetilde M_{24}$ quantum correction as order-$6$ anomaly cocycle]
\label{rem:qc-2}
The $\widetilde M_{24}$ quantum correction attaches the anomaly
cocycle of order $6$ in $H^3(M_{24}, U(1)) \supseteq \mathbb Z/12$
(generalised-Mathieu-moonshine 't~Hooft anomaly,
Gaberdiel--Persson--Ronellenfitsch--Volpato 2013) to the
$M_{24}$-equivariant quantisation of the K3 chiral structure. At the
ordered bar--cobar deformation level this is the projective
$\widetilde M_{24}$ cover induced by the $H^3$ anomaly: the
$M_{24}$-action on the K3 chiral structure is a categorical action
of the $\mathbb Z/6$-banded $2$-group classified by $[\omega] \in
H^3(M_{24}, U(1))$, restricting on the $\mathrm{Sym}_5 \subset M_{24}$
stabiliser of the umbral data to the order-$2$ Schur cocycle of
$M_{12}$ (Cheng--Duncan--Harvey~\cite{ChengDuncanHarvey2014} umbral
shift $m_\sigma$ at conjugacy class $\sigma$), detecting the obstruction
to lifting the $M_{24}$-action on the K3 elliptic genus to a strict
action on the chiral side.
\end{remark}

\begin{remark}[Classical $r$-matrix of the K3 chiral bialgebra]
\label{rem:k3-classical-rmatrix}%
\index{K3 chiral bialgebra!classical r-matrix}%
\index{Yang--Baxter equation!dynamical, Siegel-elliptic}%
\index{r-matrix!fourth class}
Assume the Vol.~III Manin-pair construction of the K3 chiral
bialgebra $\mathbf H_{\Delta_5}$
(Definition~\ref{V3-def:k3-chiral-bialgebra}).  The datum imported
into this chapter is the triple
\[
\bigl(\mathbf H_{\Delta_5},\; r^{K3}(Z),\;
K^{\kappa_{\mathrm{ch}}}=2c_+(\mathrm{Mukai}(K3))=8\bigr).
\]
Its three scalar shadows have different targets:
\[
\kappa_{\mathrm{BKM}}(\Delta_5)=5,\qquad
\operatorname{rk}(\Lambda_{K3})=24,\qquad
K^{\kappa_{\mathrm{ch}}}=8.
\]
The first is the Borcherds weight, the second is the Heisenberg
rank in the abelian residue, and the third is the Mukai conductor.
The third belongs to the $\mathcal B$-row complementarity
normalization.  Theorem~D uses the scalar $\kappa(\cA)$ in its own
target, separately from all three K3 shadows.

The classical spectral datum is outside the rational affine
Kac--Moody kernel of \eqref{eq:qc-trace-kz-normalizations}.  The home
Lie algebra is the Manin pair
$(\mathfrak{g}_{\Delta_5},\, \mathfrak{n}_{+}^{\mathrm{imag}}\oplus
\mathfrak{h}^{\mathrm{imag,rk\,23}})$ of the Gritsenko--Nikulin
Borcherds--Kac--Moody superalgebra with Cartan form on
$\Lambda^{3,2}\otimes\mathbb{R}$ of signature $(3,2)$, and the spectral
parameter is the $\mathbb{H}_2$-period matrix
$Z = \bigl(\begin{smallmatrix}\tau & z\\ z & \rho\end{smallmatrix}\bigr)$.
The classical $r$-matrix is the Siegel Kronecker--Eisenstein series
used in Vol.~III Section~\ref{V3-sec:R-Sieg-dynamical}:
\[
 r^{K3}(Z)
 \;=\;
 \sum_{\alpha\in \Delta^{\mathrm{im}}_+(\mathfrak{g}_{\Delta_5})}
 F^{\mathrm{Sieg}}_z(\langle\alpha,Z\rangle)\, e_\alpha\otimes e_{-\alpha}
 \;+\;
 \sum_{i,j\in \mathfrak{h}^{\mathrm{imag,rk\,23}}}
 \Omega_{ij}\, h_i\otimes h_j,
\]
Here $F^{\mathrm{Sieg}}_z$ is the Siegel extension of the
Kronecker series on $\mathbb{H}_2$ from Vol.~III and $\Omega$ is the
BKM Cartan form on $\Lambda^{3,2}$.  This is
\emph{Siegel-elliptic dynamical} in the sense of
Felder~\cite{Felder94} extended to
$\mathrm{Sp}_4$.  It lies outside the three one-parameter kernels
\[
r^{\mathrm{coll}}(z)=\frac{k\Omega_{\mathrm{tr}}}{z},\qquad
r^{U_q}(z)=k\Omega\,\frac{z+1}{z-1},\qquad
r^{\mathrm{Felder}}(u,\tau).
\]
CYBE holds on the paramodular locus
$K(1)\subset \mathbb{H}_2/\mathrm{Sp}_4(\mathbb{Z})$ by the imported
Vol.~III cyclic Siegel--Kronecker identity
\[
 F^{\mathrm{Sieg}}_z(u_{12}+\lambda,Z) + F^{\mathrm{Sieg}}_z(u_{23}+\lambda,Z)
 + F^{\mathrm{Sieg}}_z(u_{31}+\lambda,Z) \;=\; 0
\]
from Section~\ref{V3-sec:R-Sieg-dynamical}.  Off
$K(1)$, on the Humbert locus $H_1\cup H_4$, CYBE acquires an elliptic
anomaly of order $\eta^{24}(\tau)\cdot\phi_{10,1}(\rho)$ whose order is
$K^{\kappa_{\mathrm{ch}}} = 8$ on $H_1$ and $16$ on $H_4$.

The abelian residue map
\[
\operatorname{Res}_{\mathrm{ab}}\colon
\operatorname{gr}_{\mathrm{ab}}\mathbf H_{\Delta_5}
\longrightarrow \mathcal H_{\operatorname{rk}(\Lambda_{K3})}
\]
sends the diagonal residue of $r^{K3}$ to the Heisenberg kernel
\[
\operatorname{Res}_{\mathrm{ab}}\bigl(r^{K3}(Z)\bigr)
=r^{\mathrm{Heis}}_{24}(z)=\frac{24}{z}.
\]
In the Vol.~III presentation the non-abelian BKM bracket enters
through Frenkel--Kac vertex closure on the K3 Fock module rather than
through the abelian Lie stratum.  If
the $\mathcal B$-family normalization
$K^{\kappa_{\mathrm{ch}}}=2c_+(\mathrm{Mukai}(K3))=8$ is imposed, the
deformation parameter satisfies
\[
\hbar^2 K^{\kappa_{\mathrm{ch}}}=-1,
\qquad
\hbar^2=-\frac18 .
\]
This normalization is external to the scalar genus-curvature
calculation in this chapter.
\end{remark}
